\section*{第 \ref{ch:b3:astrophysics} 章——天体物理学}

\begin{solution}{exo:b3:astrophysics:1}
(a) $P_{\text{c}} \sim GM^2/R^4 \approx \qty{1.1e15}{Pa}$——相当于
一百亿个大气压。(b) $\bar\rho \approx \qty{1400}{kg/m^3}$：
$T \sim P m_{\text{p}}/\rho k_{\text{B}} \approx \qty{e8}{K}$ ---
粗略，高出真实值 \qty{1.6e7}{K} 一个数量级，但量级
正确。(c) 恰好处于隧道效应的阈值。(d) 原恒星
收缩并升温，\emph{直至}核心点燃，然后
停止收缩：任何质量足够大的气体球都必然抵达点燃条件
——恒星自行调节至核热源。
\end{solution}

\begin{solution}{exo:b3:astrophysics:2}
(a) $L = 4\pi d^2\times1360 \approx \qty{3.8e26}{W}$。(b) $T =
(L/4\pi R^2\sigma)^{1/4} \approx \qty{5800}{K}$。(c) 维恩定律：
$\lambda_{\text{峰}} \approx \qty{500}{nm}$——人眼进化至
太阳的峰值处。(d) 平方反比通量（本科一年级）；
斯特藩–玻尔兹曼与维恩：\cref{ch:b3:photons-phonons}中的
黑体辐射，由光子统计导出。
\end{solution}

\begin{solution}{exo:b3:astrophysics:3}
(a) 天狼星 B：远在左下；参宿四：远在右上。
(b) $R/R_\odot = \sqrt{L/L_\odot}\,(T_\odot/T)^2$：天狼星 B
$\approx 0.009\,R_\odot \approx$ 一个地球半径；参宿四
$\approx 860\,R_\odot \approx \qty{4}{au}$。(c) 水星、金星、
地球和火星都将\emph{在其内侧}运行。(d) 尺度与地球相当，
质量却与太阳相当：只有电子简并压强
（\cref{ch:b3:quantum-statistics}）——天狼星 B 是典型的
白矮星。
\end{solution}

\begin{solution}{exo:b3:astrophysics:4}
(a) $10\,\text{Gyr}\times10^{-2.5} \approx \qty{30}{Myr}$。(b)
$0.5^{-2.5} \approx 5.7$：约 570 亿年——每颗曾经诞生的红矮星
仍在燃烧。(c) 光度随
$M^{3.5}$ 增长，而燃料库只随 $M$ 增长：富者烧光其家财
比积攒还快。(d) 大质量恒星在诞生后几
百万年内死亡——尚未漂移到别处：超新星
在孕育它们的恒星形成旋臂内爆发。
\end{solution}

\begin{solution}{exo:b3:astrophysics:5}
(a) $E = -K$：辐射使 $E$ 变得更负，因此 $K$——
即温度——\emph{升高}：引力的热力学
特征。(b) $GM^2/R \approx \qty{3.8e41}{J}$。(c)
$E/L_\odot \approx \qty{e15}{s} \approx \qty{30}{Myr}$。(d)
开尔文认为太阳的年龄不可能超过数千万年，而
达尔文的蚯蚓与赖尔的岩层却要求数十亿年：
答案是当时无人知晓的一种燃料——\cref{ch:b3:nuclear-physics}，
比其早了七十年。
\end{solution}

\begin{solution}{exo:b3:astrophysics:6}
(a) $\bar\rho \approx \qty{1.8e9}{kg/m^3}$：一茶匙
的量重达九吨。(b) $g = GM/R^2 \approx \qty{3.3e6}{m/s^2}$——相当于
三十万个地球重力。(c) $n \propto M/R^3$：
简并压提供 $P \propto M^{5/3}/R^5$，引力要求
$P \propto M^2/R^4$；令其相等，得 $R \propto M^{-1/3}$。(d) 质量越大
需要更大的费米压强，而只有压缩才能提供
提供——因此质量\emph{使恒星收缩}；这一推理链有一道
极限：一旦电子进入相对论性区，压强律便减弱为
$n^{4/3}$，无论何种半径，平衡都将失效，而
$M_{\text{Ch}} = 1.4\,M_\odot$ 即为悬崖边缘。
\end{solution}

\begin{solution}{exo:b3:astrophysics:7}
(a) $V = M/\rho_{\text{核}} \approx \qty{1.2e13}{m^3}$：
$R \approx \qty{14}{km}$——一个有天际线的原子核。(b)
$\omega \propto 1/R^2$：自旋加快为 $(10^8/1.4\times10^4)^2
\approx 5\times10^{7}$，因此 25 天的周期缩短为
${\sim}\qty{40}{ms}$：进入脉冲星领域。(c) $g \approx
\qty{e12}{m/s^2}$：此处的一座山，在那里仅有毫米量级。
(d) $B \approx 10^{-2}\times5\times10^{7} \approx
\qty{5e5}{T}$——磁轴
与自转轴错位，使射电辐射沿其两极发出：光束
如灯塔般扫掠，这些闪光便被称为脉冲星。
\end{solution}

\begin{solution}{exo:b3:astrophysics:8}
(a) 由 $\tfrac12mc^2 = GMm/R$ 得 $R = 2GM/c^2$。(b) 太阳：
\qty{3.0}{km}；地球：\qty{8.9}{mm}——无论将二者中哪一个压缩到
此半径之内，光便无法逃逸。(c) $R_{\text{s}} \approx
\qty{1.2e10}{m}$：约为水星轨道的五分之一——且已观测到恒星
绕其高速旋转。(d) $\bar\rho \propto
M/R_{\text{s}}^3 \propto 1/M^2$：一个 $10^8\,M_\odot$ 的黑洞比
水还稀疏。视界处无需有任何东西“被压碎”——
它是一个不归点，而非岩石表面。
\end{solution}

\begin{solution}{exo:b3:astrophysics:9}
(a) $H_0 = 7\times10^4/3.09\times10^{22} \approx
\qty{2.3e-18}{s^{-1}}$。(b) $1/H_0 \approx \qty{4.4e17}{s}
\approx 14$ 十亿年。(c) $v = \qty{1.4e4}{km/s}$；
$z \approx v/c \approx 0.047$：波长到达时被拉长
5\,\%。(d) 膨胀速率并非恒定——早期被物质减速，
此后则由暗能量加速：这一
$1/H_0$ 仍与真实的 13.8 十亿年相差不足一个十亿年，这是
一桩小小的宇宙巧合。
\end{solution}

\begin{solution}{exo:b3:astrophysics:10}
(a) $\lambda_{\text{峰}} = b/T \approx \qty{1.1}{mm}$：
微波，正如所预期。(b) 波长被拉伸了 $3000/2.725 \approx
1100$——此后宇宙在每个方向上
都扩大了一千一百倍。复合发生在 \qty{3000}{K}，远
低于 $\qty{13.6}{eV}/k_{\text{B}}$，因为萨哈平衡
（\cref{prop:b3:grand-canonical:saha}）被每个重子对应的二十亿个光子
所偏移：电离的尾部使氢保持解离，
比朴素的能量判据所给出的时间长得多。(c) 约
每个重子约对应 $1.6\times10^{9}$ 个光子——按数目计，宇宙
几乎完全由光构成。(d) 宇宙微波背景不是某个地方，而是
整个天空：每条视线都止于最后散射的
“雾”，因此任何天线无论指向何方都能接收到它。
\end{solution}

\begin{solution}{exo:b3:astrophysics:11}
(a) $n/p = \eu^{-1.29/0.8} \approx 0.20$。(b) 取 $n/p = 1/7$：
每 2 个中子对应 14 个质子；这 2 个中子与 2 个质子结合成
一个 $^4$He（质量数为 4），余下 12 个孤立质子：$4/16 = 25\,\%$（按
质量计）。(c) 没有恒星来得及把全部物质的四分之一
变成氦；在最古老、金属量最低的气体中找到这一下限，
正是炽热的早期宇宙被当场抓获：天空中写着一个
玻尔兹曼因子。(d) $A = 5$ 或 $8$ 处不存在稳定的
原子核，而密度与温度在数分钟之内
迅速下降：通向碳的桥梁（三个氦核同时
碰撞）只在红巨星稠密而耐心的核心中才会打开。
\end{solution}

\begin{solution}{exo:b3:astrophysics:12}
(a) $mv^2/r = GM(r)m/r^2$。(b) $M = v^2r/G \approx
\qty{1.1e42}{kg} \approx 5\times10^{11}\,M_\odot$。(c) 约
十分之一：九成的星系既不发光，也不
吸收。(d) $v$ 平坦意味着 $M(r) \propto r$——不可见的
质量在可见盘外仍不断累积，形成
一个延展的暗晕（密度 $\sim1/r^2$），其中发光的
星系只是被照亮的核心。
\end{solution}

\begin{solution}{pb:b3:astrophysics:1}
\textbf{1.} $t \sim 10\,\text{Gyr}\times20^{-2.5} \approx
\qty{6}{Myr}$：太阳将活过自身寿命的两千倍。
\textbf{2.} 较重的原子核带更多电荷：更高的库仑
势垒需要更热的核心，才能使隧道效应的几率
（见\cref{exo:b3:nuclear-physics:8}）跟上。
\textbf{3.} 铁位于 $B/A$ 曲线之巅：使其聚变
\emph{耗费}能量——恒星这座熔炉的燃料阶梯止于灰烬。
\textbf{4.} 钱德拉塞卡质量：电子简并压所能承载的
最大质量；核心是一颗在活恒星内部长成的白矮星，
在 $1.4\,M_\odot$ 时其支撑便告失效。
\textbf{5.} 氢：数百万年；氦：数十万年；
硅：一天——每个阶段每千克释放的能量更少，而
恒星愈发炽热，（主要以中微子的形式）消耗得愈来愈快：
时钟在对数意义上疯狂地走。
\textbf{6.} $t \sim 1/\sqrt{G\rho} =
1/\sqrt{6.7\times10^{-11}\times10^{12}} \approx \qty{0.1}{s}$：
核心在十分之一秒内便从恒星底下塌陷下去。
\textbf{7.} $E \sim GM^2/R \approx \qty{4e46}{J}$。
\textbf{8.} 光子辐射能量为 $10^9\times3.8\times10^{26}
\times2.6\times10^{6} \approx \qty{e42}{J}$——占其中的万分之
一。其余能量以中微子形式离开：坍缩的核物质
对其他一切都不透明，但
（见\cref{ch:b3:particle-physics}）几乎对弱相互作用
自身的粒子透明，它们在数秒之内便带走了能量。
\textbf{9.} $E/4\pi d^2 = 4\times10^{46}/2.8\times10^{43}
\approx \qty{1.4e3}{J/m^2}$。
\textbf{10.} 按每个 $\qty{15}{MeV} = \qty{2.4e-12}{J}$ 计：
${\sim}6\times10^{14}$ 个中微子每平方米——穿过
地球的每一平方米。
\textbf{11.} $6\times10^{14}\times10^{-46}\times1.5\times
10^{32} \approx 9$：探测器记录了十一个——一次
数量级上的正中靶心：这颗恒星死于
160\,000 年前。
\textbf{12.} 中微子在坍缩时即离开核心；而
光只有在激波冲破恒星的
包层、数小时之后才产生。两者历经 160 千年飞行后
同夜抵达，由此将中微子速度限定在与光速相差
$10^{9}$ 分之几以内——并使其质量几乎为零。
\textbf{13.} 能量预算（$\sim10^{46}\,\unit{J}$，99\,\%
以中微子形式存在），以及 ${\sim}10$ 秒的爆发时长——
二者都与理论中的坍缩核心相吻合，使超新星机制
从黑板上的理论走向观测。
\textbf{14.} $\bar\rho = 2.8\times10^{30}/(\tfrac43\pi
(1.2\times10^4)^3) \approx \qty{4e17}{kg/m^3}$：核
密度——\cref{def:b3:nuclear-physics:nuclide}中的茶匙，
如今却有一座城市那么大。
\textbf{15.} 自转加快 $(10^8/1.2\times10^4)^2 \approx
7\times10^{7}$：30 天 $\to$ ${\sim}\qty{40}{ms}$。
\textbf{16.} $B \approx 10^{-2}\times7\times10^{7} \approx
\qty{7e5}{T}$——一亿特斯拉·米的磁通量，
守恒地压缩进一张邮票大小的面积里。
\textbf{17.} 磁轴偏离自转轴，
沿磁极发射射电束；这束射电像
灯塔一样扫过天空。蟹状星云的遗迹——1054 年的客星——仍
每秒自转三十圈，千年的自转减慢才刚刚
开始。
\textbf{18.} $R_{\text{s}} = 3\times\qty{3}{km} \approx
\qty{9}{km}$——与恒星本身同一量级：再无
压力与之抗衡，坍缩将穿过视界继续进行，
遗迹便是黑洞。
\textbf{19.} 这类时钟阵列是引力波
探测器，其孔径可达星系尺度（也是相对论的最佳
实验室：双脉冲星轨道的衰减恰与
引力辐射的预言一致）。
\textbf{20.} 其中的氧在大质量恒星的壳层中聚变生成
并由超新星抛射出来；铁则在
爆炸的终点本身锻造而成；二者漂入星际云，
随后汇入太阳星云、地球、海洋与食物链
——乃至读者本人。不过大爆炸的氢，
这一份，一直都是你的。
\textbf{21.} 爆炸质量相同 $\to$ 固有亮度相同
$\to$ 视亮度给出距离：一支
瓦数已知的蜡烛，在数十亿光年外仍可见——
这把标尺画出了哈勃图的远端。
\textbf{22.} 宇宙膨胀正在\emph{加速}：某种
“暗能量”将宇宙推开——2011 年的诺贝尔奖，
也是物理学最大的一笔未结账单。
\textbf{23.} $\qty{50}{kpc} \approx 163000$ 光年：
这束光离开之时，\emph{智人}刚刚开始走出非洲。
\textbf{24.} $6\times10^{14}\times0.03 \approx 2\times10^{13}$
穿过每一个在世的人——而且，按
\cref{exo:b3:particle-physics:10}给出的概率，百万人中
竟无一人与哪怕一个中微子发生过相互作用：那道比整个星系
更亮的闪光，无声地穿过了人类。
\textbf{25.} 一颗恒星依靠隧穿质子维持了千万年的平衡；
十分之一秒的坍缩以中微子形式付出了 $10^{46}$
焦耳；泡利原理将残骸
维持在核密度；而散落的残片——氧、铁、
读者——仍乘哈勃膨胀而运动，在一个
婴儿照片是普朗克曲线的宇宙中。物理学暂告完备：
本分册到此结束。
\end{solution}
